\section{Very-general jet estimates}\label{proj:sec:jets}

We retain the counterexample in Equation~\eqref{proj:ex:counterexample}
and the induction hypothesis of Assumption~\ref{proj:mmp:lower-models}.
Thus \(X\) is smooth projective,
\(K_X\) is pseudo-effective, and \(\kappa(X,K_X)=-\infty\), whereas
good minimal models exist in smaller dimensions.  We use
Propositions~\ref{proj:prop:general-type}, \ref{proj:prop:bounded-curves}, and
\ref{proj:prop:signed-alternative}.  Dimension one is impossible, since a
smooth curve with pseudo-effective canonical divisor has genus at least
one.  Hence \(n=\dim X\ge2\).

Fix the late terminal \(\Q\)-factorial model \(Y\) from
Proposition~\ref{proj:prop:bounded-curves}.  Write \(K=K_Y\) and \(A=A_Y\).
The subsequent scaling models \(Y_t\) are small modifications of \(Y\);
the transforms \(A_t\) are effective up to \(\Q\)-linear equivalence,
and \(K_t+tA_t\) is nef, big, and semiample.  Put
\[
 \mu_t=\vol(X,K_X+tA)^{1/n}.
\]
Here, in the volume on \(X\), \(A\) denotes the original ample divisor.
We have \(\mu_t\to0\).  Choose an integer \(m>0\) with \(mK\) Cartier.

\subsection{Normalization and two local tools}

Fix a small rational number \(\epsilon>0\).  An integer \(q>0\) will
be chosen after \(\epsilon\) and before \(t\).  As rational \(t\downarrow0\),
choose positive integers \(r=r(t)\) with
\begin{equation}\label{proj:jet:normalization}
 r\mu_t\longrightarrow\epsilon,\qquad
 L=r(K+tA),\qquad L_b=L+bmK\quad(0\le b\le q).
\end{equation}
Only rational weights \(b\) are used for model constructions;
volume functions in integrals are extended continuously to real weights.
The positive part of \(L_b\) is
the nef semiample class
\[
 N_b=(r+bm)(K_s+sA_s)
 \quad\hbox{on }Y_s,\qquad s=\frac{rt}{r+bm}.
\]
When emphasizing the weight, denote this axis model by \(Y(b)=Y_s\)
and write \(A_b\) for its transform of \(A\).
All references to positive parts below mean these classes on their
scaling models, or their pullbacks to common resolutions.  Monotonicity
of volume in pseudo-effective order gives
\begin{equation}\label{proj:jet:volume-comparison}
 \vol(L)\le \vol(L_b)\le
 \left(1+\frac{qm}{r}\right)^n\vol(L).
\end{equation}
Indeed \(L_b-L=bmK\) is pseudo-effective, and
\((1+bm/r)L-L_b=bmtA\) is effective up to equivalence.
Consequently \(N_b^n\) is bounded above and bounded away from zero,
uniformly for \(b\in[0,q]\), and tends uniformly to \(\epsilon^n\).
More explicitly, for every fixed \(q\), once \(t\) is sufficiently
small depending on \(q\),
\begin{equation}\label{proj:jet:uniform-normalized-volume}
 \frac{\epsilon^n}{2}\le N_b^n\le2\epsilon^n
 \qquad(0\le b\le q).
\end{equation}
The displayed constants are independent of \(q\).
All unspecified constants in this section may depend on
\(n,\epsilon,q,Y,A,m\), but not on sufficiently small \(t\) or on \(b\).

A family of curves with uniformly bounded geometric genus and
\(N_b\)-degree cannot cover the corresponding \(Y_s\) for arbitrarily
small \(t\).  In fact \(s/t\to1\) uniformly in \(b\), and
\[
 (r+bm)\mu_s=\vol(L_b)^{1/n}
\]
is bounded above and away from zero.  Thus such curves would have
uniformly bounded \((K_s+sA_s)\)-degree divided by \(\mu_s\), contrary
to Proposition~\ref{proj:prop:bounded-curves}.  We call this consequence the
\emph{normalized curve exclusion}.

The following numerical form of subadjunction is useful because an
auxiliary pair need only be lc at the generic point of its center.
All intersections with a divisor on a possibly nonnormal center mean
intersections after normalization and resolution.

\begin{lemma}[Numerical subadjunction]\label{proj:jet:subadjunction}
Let \(Z\) be projective and \(\Q\)-factorial klt, let
\(\Theta\ge0\) be rational, and let \(V\) be an lc center of
\((Z,\Theta)\) at whose generic point the pair is lc.
If \(f=\dim V>0\), \(P\) is a nef \(\Q\)-Cartier class on \(V\), and
\(V^*\to V\) is a resolution, then
\begin{equation}\label{proj:jet:subadjunction-bound}
 K_{V^*}P^{f-1}\le (K_Z+\Theta)|_V P^{f-1}.
\end{equation}
\end{lemma}

\begin{proof}
The dlt modification for an arbitrary effective boundary
\cite[Theorem~2.10]{Proj-FujinoHashizume2023}, applied after truncating
coefficients exceeding one, gives
\[
 K_Q+B+E=p^*(K_Z+\Theta),
\]
where \(Q\) is \(\Q\)-factorial, \((Q,B)\) is dlt, and \(E\ge0\)
is supported above the non-lc locus.  Choose, over the generic point
of \(V\), a minimal lc stratum, and denote its closure by \(S\).
The divisor \(E\) does not contain \(S\).  Iterated dlt adjunction,
including the effective restriction of \(E\), produces an effective
divisor \(B_S\) with
\[
 K_S+B_S\sim_{\Q}
 (p|_S)^*((K_Z+\Theta)|_V).
\]
The pair is klt over the generic point of \(V\); singularities elsewhere
are not asserted to be klt.

Factor \(S\to V^\nu\) through its Stein base \(V'\).  This is a
klt-trivial fibration: generic klt, effectivity of the boundary, and
connected fibers give the rank-one condition.  The canonical bundle
formula, in precisely this generically subklt form, gives a
discriminant and a b-nef moduli b-divisor
\cite[Definition~3.1 and Theorem~3.5]{Proj-FujinoGongyoModuli2014}.
The discriminant trace on \(V'\) is effective.  To see the sign,
over a generic base prime a divisor in its inverse image has
multiplicity at least one and a nonnegative boundary coefficient;
the lc threshold of the fiber is therefore at most one.
Coefficients exceeding one in the discriminant cause no difficulty
for the present inequality.

Intersect the resulting base formula with the pullback of \(P^{f-1}\).
The moduli term is nonnegative: on a model where it is nef this is a
nef intersection, and exceptional terms vanish on pushing down.
The effective discriminant is also nonnegative.  Finally the
codimension-one ramification formula for the finite map
\(V'\to V^\nu\), followed by projection, can only increase the
canonical intersection.  Discrepancy terms on resolutions have
images of codimension at least two and pair trivially with the
pulled-back \(P^{f-1}\).  Dividing by the finite degree gives
Equation~\eqref{proj:jet:subadjunction-bound}.
\end{proof}

\begin{lemma}[Tracking a moving base component]\label{proj:jet:tracking}
Let \(P\) be nef, semiample, and big on a projective \(d\)-fold \(Z\).
Work at very general smooth marked points \(x\).  Choose \(d+1\)
levels
\[
 \tau_0<\tau_1<\cdots<\tau_d,\qquad
 \tau_{i+1}-\tau_i=\delta>0.
\]
Suppose, in sufficiently large divisible degree \(k\), all the spaces
\[
 {\cal V}_{i,x}
 =H^0\bigl(Z,kP\otimes\mathfrak m_x^{\lceil k\tau_i\rceil}\bigr)
\]
are nonzero, and the base locus of \({\cal V}_{0,x}\) has a
positive-dimensional component through \(x\).
Then one obtains an algebraic family of integral components \(V=V_x\)
sweeping \(Z\), of dimension \(0<f<d\), and a step for which \(V\)
is a component of both successive base loci.  At its generic point
the higher subseries has an isolated base component \(V\) and
vanishes to order at least \(k\delta/2\), once \(k\) is sufficiently
large.  Moreover
\begin{equation}\label{proj:jet:tracking-estimates}
 \begin{split}
  P^f.V&\le (2/\delta)^{d-f}P^d,\\
  K_{V^*}P^{f-1}
   &\le (K_Z+cP)|_V P^{f-1},
       \qquad 0<c\le 2d/\delta,
 \end{split}
\end{equation}
whenever \(Z\) is \(\Q\)-factorial klt.
\end{lemma}

\begin{proof}
The base loci increase with the level.  Following containing
irreducible components through the marked point produces a nested
chain of proper positive-dimensional subvarieties.  There are only
\(d-1\) possible dimensions, so two successive components agree.
For fixed degree, jet kernels form vector bundles after shrinking
the marked-point parameter space.  Components of their base loci
spread after a finite parameter extension and further shrinking.
We may fix the chosen step, dimension, and component on an
irreducible parameter space \(T\).  Its incidence
\(\Gamma\subset T\times Z\) dominates \(Z\), because it contains the
varying marked point.

Here is the multiplicity argument.  Regard a local section of the
higher kernel bundle as a varying element of the fixed vector space
\(H^0(Z,kP)\).  A parameter derivative of order \(j\) loses at most
\(j\) orders of vanishing along the moving diagonal.  Thus, for
\(j<k\delta-O(1)\), every such derivative belongs to the lower kernel
and vanishes on the selected incidence \(\Gamma\).
At a general smooth point of \(\Gamma\), the projection
\(\Gamma\to Z\) is smooth.  Therefore the pure parameter directions,
together with \(T\Gamma\), span \(T(T\times Z)\).
In local coordinates over \(Z\), this says that normal coordinates
to \(\Gamma\) can be chosen among parameter coordinates.  Vanishing
of all parameter derivatives of orders below \(h\) is consequently
equivalent there to membership in \({\cal I}_\Gamma^h\).
Restricting to a parameter fiber proves generic multiplicity at
least \(k\delta-O(1)\) along \(V\), hence at least \(k\delta/2\).
This applies to a local basis of the higher kernel bundle, and
therefore to every section of the higher subseries.

Put \(a=d-f\) and \(h=\lceil k\delta/2\rceil\); the stronger bound
\(k\delta-O(1)\) permits this choice for sufficiently large \(k\).
At the smooth generic point of
\(V\), the base ideal is maximal-ideal primary and lies in the
\(h\)-th power of that maximal ideal.  Cut by \(a\) general members
of the subseries.  Their local intersection multiplicity along
\(V\) is at least \(h^a\).  For completeness, global excess base
components do not invalidate this bound: at each cut discard
components wholly contained in the base locus before the next
intersection.  None contains the generic point of \(V\), since
\(V\) is a base-locus component.  The discarded cycles have
nonnegative \(P\)-degree by nefness.  The remaining proper
intersections therefore have total \(P^f\)-degree at most
\(k^aP^d\).  Since \(h\ge k\delta/2\), this gives the first inequality.

The local lc threshold of this primary ideal is at most \(a/h\):
the exceptional divisor of the blowup of its smooth closed point
in the \(a\)-dimensional transverse local scheme gives this bound.
On a log resolution, a sufficiently long average of general
members realizes the ideal threshold by an effective rational
divisor \(\Theta\sim_{\Q}cP\).  It is lc at the generic point of
\(V\), with \(V\) an lc center, and
\(c\le k a/h\le2d/\delta\).
Lemma~\ref{proj:jet:subadjunction} gives the second inequality.
\end{proof}

We record how these estimates give curves rather than merely
bounded intersection numbers.  Suppose a moving \(f\)-fold \(V\)
is of general type, \(P|_V\) is nef and big, and
\begin{equation}\label{proj:jet:curve-input}
 0<P^f.V\le C_0,\qquad
 K_{V^*}P^{f-1}\le C_1P^f.V.
\end{equation}
Effective birationality gives a uniform pluricanonical degree whose
moving part, on a further resolution, is a big basepoint-free
Cartier divisor \(H\) defining a birational morphism
\cite[Theorem~4.0.1]{Proj-HMX2018}.  Thus
\(HP^{f-1}\le C_2P^f.V\).  Mixed Hodge index gives
\[
 H^jP^{f-j}\le C_2^jP^f.V\quad(0\le j\le f);
\]
no lower bound on \(P^f.V\) is needed here.  For \(f>1\), cut by
\(f-1\) general members of \(|H|\).  The resulting integral curves
have bounded \(P\)-degree and are birational to linear curve
sections of a projective variety of bounded degree.  Generic plane
projection bounds their geometric genus.  For \(f=1\), the
canonical-degree bound already bounds the genus.
The same argument works for a log smooth pair of log general type
with reduced boundary, using coefficients in the fixed set
\(\{1\}\) in effective birationality.

If such components sweep an ambient space mapping to \(Y\), and
their curves project nonconstantly with \(N_b\)-degree bounded by
their \(P\)-degree, we obtain forbidden covering curves on \(Y_s\).
Nonconstant projection follows by choosing the complete intersections
generally for the big birational system.  The curves can be
parameterized in covering algebraic families: the components sweep,
the choices can avoid any prescribed countable exceptional union,
and Hilbert schemes and spaces of maps have only countably many
components.  We will use this consequence of
Equation~\eqref{proj:jet:curve-input} repeatedly.

\begin{proposition}[Scalar jet bound]\label{proj:prop:scalar-jets}
For sufficiently small \(t\), let \(P\) be the positive part of
\(2r(K+2tA)\), and put \(\rho=(P^n)^{1/n}\).
At a very general point, every nonzero section of every divisible
multiple \(kP\) has order at most \(4k\rho\).
The same statement holds on a common resolution after adding an
effective exceptional divisor that does not change the section
spaces.  Moreover
\[
 2r\mu_t\le\rho\le4r\mu_t,
\]
so in particular \(\rho\le8\epsilon\) for small \(t\).
\end{proposition}

\begin{proof}
The volume bounds follow from pseudo-effectivity of \(K\):
\[
 \vol(K+tA)\le\vol(K+2tA)\le2^n\vol(K+tA).
\]
Suppose the asserted order bound fails for arbitrarily small \(t\).
Taking powers of the offending section permits arbitrarily large
divisible degrees.  Choose \(n+1\) levels strictly between \(\rho\)
and \(4\rho\), with equal gaps comparable to \(\rho\).
All corresponding subseries are nonzero.  The base point at the
lowest level cannot be isolated: local Bezout for \(n\) general
members would exceed the total intersection \(k^nP^n\).
Lemma~\ref{proj:jet:tracking} therefore supplies a moving proper
positive-dimensional component \(V\).

The component is of general type by
Proposition~\ref{proj:prop:general-type}.  On the scaling model for
\(K+2tA\), the effective transform of \(A\) does not contain a
general such component, and
\(K_Z|_V\le(2r)^{-1}P|_V\) in effective order.
Equations~\eqref{proj:jet:tracking-estimates} consequently imply
Equation~\eqref{proj:jet:curve-input} with constants uniform in \(t\);
the gap is bounded above and away from zero because
\(\rho\) is comparable to the fixed \(\epsilon\).
The resulting bounded-genus, bounded-normalized-degree covering
curves contradict Proposition~\ref{proj:prop:bounded-curves}.
Exceptional additions preserve sections and their orders at
general points, proving the additional assertion.
\end{proof}

\subsection{The two-slot bundle}

On \(Y\times Y\), let
\begin{equation}\label{proj:jet:two-slot}
 Z_0=\PP(mK_1\oplus mK_2),\qquad
 \xi=\OO_{Z_0}(1),\qquad M=L_1+L_2+q\xi.
\end{equation}
We use the convention that sections of \(k\xi\) are symmetric powers
of the indicated direct sum.  Fixing one base slot gives the
one-slot bundle
\(\PP(\OO_Y\oplus mK)\), with class \(L+q\xi\), up to a constant
line from the fixed slot.  We call it a \emph{slice}.
The torus open in either bundle is the complement of its two axes.

\begin{lemma}[Models, volume, and weights]\label{proj:jet:bundle-models}
Both the total class \(M\) and the slice class have big semiample
positive parts \(P\) on \(\Q\)-factorial klt models \(Z\), with
\[
 K_Z+\Delta\sim_{\Q}c_tP,\qquad \Delta\ge0,
\]
where \(c_t\) is bounded independently of small \(t\).
Their volumes satisfy, respectively,
\begin{align}
 \vol(M)&=\frac{(2n+1)!}{(n!)^2}
   \int_0^q\vol(L_b)\vol(L_{q-b})\,db,\label{proj:jet:total-volume}\\
 \vol(L+q\xi)&=(n+1)\int_0^q\vol(L_b)\,db.\label{proj:jet:slice-volume}
\end{align}
In particular both volumes are bounded above and below by positive
constants times \(q\).  For each rational weight, on common
resolutions,
\begin{equation}\label{proj:jet:weight-domination}
 P\sim_{\Q}N_{b,1}+N_{q-b,2}+E_b,\qquad E_b\ge0,
\end{equation}
with the analogous one-term formula on slices.
At a generic point of a base divisor of a slice, and on a general
torus fiber, the full system has no fixed subtraction.
\end{lemma}

\begin{proof}
The product \(Y\times Y\) is \(\Q\)-factorial.  Indeed on a product
resolution the Picard group has no cross term because the smooth
models have irregularity zero; the two factorwise
\(\Q\)-factorial descent statements then apply.  Products of
canonical singularities are canonical, and the projective bundles
are klt.  The two disjoint Cartier axes form a plt boundary \(D\),
and
\[
 K_{Z_0}+D=K_{\rm sum},\qquad
 D\sim2\xi-mK_{\rm sum}.
\]
The notation \(K_{\rm sum}\) means \(K_1+K_2\), or \(K\) on a slice.

For \(0<\sigma\le1\), put
\[
 \gamma=2\sigma+\frac{q(1+\sigma m)}r,\qquad
 \Xi_t\sim_{\Q}\xi+
       \frac{(1+\sigma m)t}{\gamma}A_{\rm sum}.
\]
Both endpoint weight systems are big, so there are effective
rational representatives of \(\Xi_t\) containing neither axis.
We first obtain a threshold bound independent of \(\sigma\).
Restrict to \(0<t\le\sigma/(1+m)\).  Since \(\gamma\ge2\sigma\),
the coefficient
\[
 a=\frac{(1+\sigma m)t}{\gamma}
\]
lies in \((0,1]\).  Thus the numerical classes
\(\xi+aA_{\rm sum}\), and their restrictions to either fixed axis,
range over bounded segments independent of \(\sigma,q,r,t\).
On a fixed smooth resolution their effective pullbacks have
uniformly bounded degree against a fixed very ample class.
That degree bounds their multiplicity at every point, including
the coefficients of exceptional components.  Rational denominators
do not enter this estimate.  The smooth lc threshold is at least
the reciprocal of maximal multiplicity.

On the fixed klt space, write the crepant boundary on its resolution
as an SNC divisor with coefficients below one.  The minimum of
one and the positive numbers one minus its positive coefficients
bounds its log discrepancies below by a fixed positive multiple
of smooth log discrepancies.  Therefore the preceding smooth
bound gives a common threshold \(\eta>0\) on the original space.
The same reasoning on the axes gives, after decreasing \(\eta\),
the same bound for both restricted divisors.

Now choose rational \(0<\sigma<\min\{1,\eta/4\}\), independently of
\(q,t\), and then take \(t\) small enough that
\(t\le\sigma/(1+m)\) and \(q(1+\sigma m)/r<\sigma\).
It follows that \(\gamma<3\sigma<\eta\).
Inversion of adjunction with the disjoint coefficient-one axes
gives plt near them; away from them the threshold bound gives klt.
Lowering their coefficients to \(1-\sigma\) consequently shows that
\[
 (Z_0,(1-\sigma)D+\gamma\Xi_t)
\]
is klt.  This choice respects the order: first \(\epsilon\), then
\(q\), then sufficiently small \(t\) with \(r\) as in
Equation~\eqref{proj:jet:normalization}.
A direct calculation gives its adjoint class
\[
 K_{Z_0}+(1-\sigma)D+\gamma\Xi_t
   \sim_{\Q}\frac{1+\sigma m}{r}M.
\]
The big klt adjoint has a good minimal model by
\cite{BCHM}.  Pushing the boundary to this model gives the stated
\(\Delta\) and \(c_t=(1+\sigma m)/r\).

For the volumes, the weight-\(l\) summand in degree \(k\) has base
classes \(kL_{l/k}\) and \(kL_{q-l/k}\); on a slice there is one
such factor.  The section decomposition and asymptotic Riemann--Roch
give Equations~\eqref{proj:jet:total-volume}--\eqref{proj:jet:slice-volume}
as Riemann sums.  One can justify passage to the integral without
assuming uniform asymptotic Riemann--Roch: partition \([0,q]\) into
rational bins, compare weight classes in each bin by adding and
subtracting the bin width times a fixed very ample divisor
dominating both \(mK\) and \(-mK\), take divisible-degree limits,
and then shrink the bins.  Volume continuity gives the formulas.
Equation~\eqref{proj:jet:volume-comparison} supplies their uniform bounds
and, in particular, proves bigness used above.
In the range of Equation~\eqref{proj:jet:uniform-normalized-volume},
the bounds needed when choosing \(q\) are explicitly
\[
 \vol(M)\ge \frac{(2n+1)!}{4(n!)^2}\,q\epsilon^{2n},
 \qquad
 \vol(L+q\xi)\le2(n+1)q\epsilon^n.
\]
Their coefficients are independent of \(q\); the smallness threshold
for \(t\) is allowed to depend on \(q\).

Compare a complete system on a common resolution with the subsystem
at a fixed rational weight.  The latter has fixed divisor consisting
of its toric monomial and the base-model fixed differences.
Subtracting the fixed divisor of the full system leaves an effective
divisor \(E_b\); its moving class is precisely the sum of the
pullbacks of the indicated \(N_b\)'s.  This proves
Equation~\eqref{proj:jet:weight-domination}.  The same comparison in
individual sufficiently divisible degrees shows that every section
at that weight, after the full fixed subtraction, contains the
corresponding multiple of \(E_b\).

At a generic base-divisor point the small base-model maps are
isomorphisms.  The two endpoint systems are free there in divisible
degree, so together they have no common zero on the projective-line
fiber.  The full system therefore has no fixed subtraction there.
The same holds on a general fiber.  Finally, when restricting
Equation~\eqref{proj:jet:weight-domination} to a component sweeping the
torus open, choose a general component not contained in
\(\Supp E_b\).  There are only countably many rational weights, so
these effective restrictions can be required simultaneously.
\end{proof}

For a nef \(\Q\)-Cartier class \(P\) on a projective variety \(Z\)
and a smooth point \(x\in Z\), its Seshadri constant is
\[
 \varepsilon(P;x)=
 \inf_{\substack{C\subset Z\ \mathrm{integral\ curve}\\x\in C}}
       \frac{P\cdot C}{\mult_x C}.
\]

\begin{proposition}[Large two-slot Seshadri constant]\label{proj:prop:large-seshadri}
For each fixed sufficiently small \(\epsilon>0\), there is an integer
\(q>0\) such that, for all sufficiently small rational \(t>0\) and
\(r\) as in Equation~\eqref{proj:jet:normalization}, the positive part
\(P\) of Equation~\eqref{proj:jet:two-slot} has
\[
 \varepsilon(P;x)>2n+2
\]
at a very general smooth point of its torus open.
In particular, for some sufficiently divisible integer \(k>0\),
the complete system \(|kP|\) separates jets of order strictly
greater than \((2n+2)k\) at such a point.
The section spaces and these general-point jets agree on the
original bundle and on common birational resolutions.
\end{proposition}

\begin{proof}
Write \(d=2n+1\).  Choose \(q\) sufficiently large that the volume
root in Equation~\eqref{proj:jet:total-volume} admits \(d+1\) levels above
\(2n+3\), with equal gap \(\delta\) as large as needed below.
This is possible uniformly for small \(t\), since the volume is
bounded below by a positive constant times \(q\).
Assume that the Seshadri assertion fails for arbitrarily small \(t\).
Jet counts show that the systems at all these levels are nonzero
in sufficiently large divisible degree.  A curve through the marked
point with degree-to-multiplicity ratio below the lowest level is
contained in that system's base locus.  Such a curve exists from
the assumed Seshadri bound, whether or not the infimum defining
the constant is attained.  Lemma~\ref{proj:jet:tracking} supplies a
moving proper component \(V\) and both estimates
\eqref{proj:jet:tracking-estimates}.

For all moving components under consideration, the effective
boundary \(\Delta\) of Lemma~\ref{proj:jet:bundle-models} does not contain
the component.  Thus \(K_Z|_V\le c_tP|_V\) in effective order.
Whenever \(V\) is of general type, the curve construction following
Equation~\eqref{proj:jet:curve-input} and weight domination give a
contradiction.  We must also handle components that are not
generically finite over a proper base image.

\paragraph{Restriction to a slice.}
Fix a general value of one slot projection of \(V\), and suppose
the corresponding fiber component \(F\) has
\(0<\dim F<n+1\).  It is an isolated high base component for the
restricted subseries on the opposite slice.

We give the local justification.  On a common isomorphic torus
open, let \(I\) be the total high-series base ideal.  Remove the
other base components from a dense open \(V^\circ\) of \(V\).
On an ambient neighborhood \(U\) of its generic point,
\(\Supp(\OO_U/I)=V\cap U\).
After restricting \(V\) and its image to smooth opens, generic
smoothness makes the general projection fiber generically reduced.
A general fiber component \(F\) meets \(V^\circ\), and the ideal
of \(V\) restricts to the ideal of \(F\) at its generic point.
Near that point,
restriction to the slice therefore has support exactly \(F\);
the restricted ideal is maximal-ideal primary there.  The inclusion
in the \(h\)-th power of the generic ideal of \(V\) restricts to
the \(h\)-th power of the generic ideal of \(F\).
The fixed differences of the full and slice models are units on
a further common open, so dividing them out does not alter this
statement.  The restricted global sections are a subseries of
the complete slice system after trivializing the fixed-slot lines.
Thus the proof of Lemma~\ref{proj:jet:tracking} applies to \(F\), using
slice volume, with multiplicity at least \(k\delta/2\).

These \(F\)'s may be chosen in families sweeping the opposite
slice.  Indeed the incidence of the \(V\)-family dominates the
total torus.  Choose a general incidence point, not necessarily
the original marked point, and then a general fiber of its first
slot evaluation.  Dominance gives the required dominant evaluation
onto the opposite slice.  Generic flatness permits the component
choices just made.  For each fixed \(q,t\), very general choices
also avoid the exceptional sets for all rational weights.

\paragraph{Proper base images on a slice.}
If \(F\) is generically finite over a proper positive-dimensional
image \(W\) in the remaining base, then \(W\), and hence \(F\), is
of general type by Proposition~\ref{proj:prop:general-type}.
The slice estimates, subadjunction, and weight domination give
forbidden bounded curves.

Otherwise \(F\) is saturated over its base image \(W\): on the
original bundle it is the full projective-line bundle over \(W\).
Put \(w=\dim W=\dim F-1\).  The nef weight comparison gives
\begin{equation}\label{proj:jet:saturated-degree}
 qN_b^w.W\le P^{w+1}.F.
\end{equation}
Indeed \(P-\pi^*N_b\) is effective on the resolved graph, so
successive mixed nef intersections give
\(P^{w+1}.F\ge P(\pi^*N_b)^w.F\).  The latter is
\(qN_b^w.W\), since endpoint freeness gives fiber degree \(q\).
If \(W\) is a point, the lower bound \(q\) contradicts the slice
Bezout upper bound \(2(n+1)q\epsilon^n(2/\delta)^n\), after \(q\)
has been chosen to make \(\delta\) large.  In particular the
coefficient used in this choice does not depend on \(q\).

Suppose \(w>0\).  We construct a small adjoint on \(W\), since the
ruled \(F\) itself need not be of general type.  In degree \(k\),
expand every section of the high-vanishing slice subseries into
its toric weights.  At the generic point of \(W\), let \(I_l\) be
the ideal generated by the coefficients of the nonzero weight \(l\)
over all these sections.  Put \(h=\lceil k\delta/2\rceil\).
All \(I_l\) lie in \(\mathfrak m_W^h\): vanishing along the generic
torus fiber says that the coefficient of every Laurent monomial
vanishes to this order.  Their sum is
\(\mathfrak m_W\)-primary.  Otherwise its larger common zero germ,
times the generic torus fiber, would contradict isolation of \(F\).

Work in the regular local ring at the generic point of \(W\), of
dimension \(c=\codim W\).  Resolve the finitely many ideals
simultaneously and consider the rational lc polytope
\[
 u_l\ge0,\qquad
 \sum_lu_l\ord_E(I_l)\le a(E;Y,0).
\]
Include the exceptional divisor of the blowup of the closed point.
Since every \(I_l\subset\mathfrak m_W^h\), it gives
\(\sum_lu_l\le c/h\).  The polytope is compact and has a positive
rational maximum for \(\sum_lu_l\).  At a maximizing point, each
coordinate participates with positive order in an active
inequality; otherwise that coordinate could be increased.
The corresponding lc center \(C_l\) is contained in \(V(I_l)\)
and contains the generic point of \(W\).  Their intersection is
contained in \(V(\sum_l I_l)\), which is \(W\) locally.  The
intersection property for lc centers
\cite[Theorem~1.7]{Proj-KollarKovacs2010}, applied to the identity
morphism on the lc open, therefore makes \(W\) an lc center
generically.  To realize the ideals, choose an integer
\(M_0>\max_l u_l\) and \(M_0\) general coefficient divisors
\(D_{l,j}\) from each system.  Set
\(\Theta=\sum_l(u_l/M_0)\sum_{j=1}^{M_0}D_{l,j}\).
On the simultaneous resolution its fixed crepant coefficients are
at most one, and its general moving divisors meet transversely
with coefficients below one.  Thus the pair is lc near the generic
point of \(W\), and every active valuation remains an lc place.
The intersection argument therefore applies to this actual
effective rational divisor.  It also applies when some maximizing
coordinate \(u_l\) is zero: the active valuation blocking that
coordinate still has center in \(V(I_l)\).

Set
\[
 b=\frac{\sum_lu_l(l/k)}{\sum_lu_l},
 \qquad c'=k\sum_lu_l\le\frac{kc}{h}.
\]
Its class on the base is \(c'L_b\).  Transform to the small scaling
model for \(L_b\), which is unchanged at the generic point under
consideration.  There its class is \(c'N_b\).  Lemma
\ref{proj:jet:subadjunction}, effective \(A_b\), and
Equation~\eqref{proj:jet:saturated-degree} give bounded adjoint
intersection and bounded \(N_b^w\)-degree for \(W\).
Since \(c'\) is bounded by a dimension-dependent multiple of
\(1/\delta\), these bounds are uniform in \(t\).
The variety \(W\) is of general type by
Proposition~\ref{proj:prop:general-type}; normalized curve exclusion
again gives a contradiction.

\paragraph{A slice multisection.}
The remaining proper positive-dimensional slice case is
\(\dim F=n\), with \(F\) generically finite of degree \(e\) over
the whole base.  Weight comparison and the slice degree bound give
\[
 eN_b^n\le P^n.F,
\]
so \(e\) is bounded uniformly in \(t\).
The closure of \(F\) on the original bundle is neither axis.
Its intersections with the two axes push to effective integral
Weil divisors \(D_1,D_2\) on \(Y\), and the projective-bundle
relation gives
\begin{equation}\label{proj:jet:axis-signed}
 D_1-D_2\sim_{\Q} \pm emK.
\end{equation}
At the appropriate endpoint weight \(b\in\{0,q\}\), the difference
\(E_b|_F\) contains the corresponding toric-axis intersection with
coefficient \(q\).  There is no full fixed subtraction above
generic base-divisor points by Lemma~\ref{proj:jet:bundle-models}.
Intersecting with \(N_b^{n-1}\) and using mixed nef comparison yields
\[
 qN_b^{n-1}.D_i\le P^n.F.
\]
The other divisor has bounded degree for the same \(N_b\), because
of Equation~\eqref{proj:jet:axis-signed} and
\[
 0\le K_{Y(b)}N_b^{n-1}\le\frac{N_b^n}{r+bm}.
\]
On a log resolution add the reduced strict support of
\(D_1+D_2\) and all exceptional divisors to the canonical divisor.
This log canonical divisor is big by
Proposition~\ref{proj:prop:signed-alternative}.
Its intersection with the pulled-back \(N_b^{n-1}\) is bounded:
exceptionals vanish under projection, and reduced support has
degree no larger than \(D_1+D_2\).
Log effective birationality and the curve construction therefore
contradict normalized curve exclusion on the base itself.

\paragraph{Reduction to a correspondence.}
We have excluded fiber dimensions strictly between zero and
\(n+1\) over either slot.  A full \((n+1)\)-dimensional fiber over
one slot would mean that \(V\) contains the whole opposite slice
over its first image \(W\).  Its fiber dimension over the other
slot would then be \(\dim W+1\), strictly between zero and \(n+1\),
unless \(V\) were the full total space.  That is impossible.
Hence \(V\) projects generically finitely in both slots and
\(\dim V\le n\).  If either image is proper, general type and the
total-space estimates give the previous curve contradiction.
The same is true if \(V\) is of general type.
Only the case \(\dim V=n\), dominant with bounded degree over both
copies of \(Y\), remains.

\paragraph{Moving branch divisors.}
Take the dominating algebraic family of these correspondences,
resolving maps after shrinking the parameter space.  Suppose a
divisorial branch component of one projection moves.  On a
resolution \(V^*\), choose a ramified divisor \(R\) above its generic
point.  There is a rational weight \(b\) for which the effective
difference \(E_b|_{V^*}\) does not contain \(R\).
Indeed take a fixed free divisible degree of \(P\); some section
does not vanish generically on \(R\).  Monomial-weight sections
span that degree, so one of them does not vanish there.
Only finitely many weights are tested in this fixed family; their
model comparisons can be resolved simultaneously.

Let \(J\) be the branch divisor on the corresponding small axis
model.  Nef comparison on \(R\), ramification into a terminal
target, and pseudo-effectivity of its canonical divisor give
\begin{equation}\label{proj:jet:branch-degree}
 N_b^{n-1}.J\le P^{n-1}.R
 \le K_{V^*}P^{n-1}\le C_q.
\end{equation}
The ramification coefficient of \(R\) is a positive integer, so is
at least one; the other ramification and exceptional terms are
nonnegative.  The last bound is the total-space subadjunction
estimate.

We spell out the adjunction needed for \(J\), without assuming
\((Y(b),J)\) globally lc.  Terminal \(Y(b)\) is smooth in codimension
two.  Thus normalization and conductor adjunction along \(J\)
give
\[
 K_{J^\nu}+C=(K_{Y(b)}+J)|_{J^\nu},\qquad C\ge0
\]
in codimension one.  On resolving \(J^\nu\), exceptional divisors
map to codimension at least two and pair trivially with the
pulled-back \(N_b^{n-2}\).  Consequently, with
\(d_J=N_b^{n-1}.J\),
\begin{align}
 K_{J^*}N_b^{n-2}
 &\le (K_{Y(b)}+J)J N_b^{n-2}\notag\\
 &\le \frac{d_J}{r+bm}+\frac{d_J^2}{N_b^n}
 \le C'_q d_J.\label{proj:jet:branch-adjunction}
\end{align}
For the second line, a general moving \(J\) is not a component of
a fixed effective representative of \(A_b\), giving the first
term; mixed Hodge index gives the second.  The lower bound for
\(N_b^n\) and Equation~\eqref{proj:jet:branch-degree} give the final
uniform constant.

A moving \(J\) sweeps very general base points and is of general
type by Proposition~\ref{proj:prop:general-type}.
Equations~\eqref{proj:jet:branch-degree}--\eqref{proj:jet:branch-adjunction}
therefore give forbidden bounded curves.  Branch components can
be named after a finite parameter extension and shrinking.
It follows that, for all sufficiently small \(t\), all divisorial
branch components in the chosen family are fixed.

\paragraph{Fixed covers.}
For each such fixed \(t\), remove the fixed branch divisors and
the singular locus from \(Y\).  Normalization and purity identify
the covers over the resulting smooth open with finite etale
covers.  Their degrees are bounded.  The topological fundamental
group of a smooth complex quasi-projective variety is finitely
generated, so it has only finitely many subgroups of bounded
index.  There are therefore only finitely many possible covers
in each slot, up to isomorphism; their finite normalizations over
\(Y\) are determined by these restrictions.

Some algebraic family of graphs of birational isomorphisms between
two fixed covers must still dominate the base product.  To justify
the family assertion, parameterize the graphs by their Hilbert
schemes, shrink for flatness and birationality of the two
projections, and use countability of these parameter spaces
together with the finite cover choices.  Identify the two covers
by one such birational isomorphism.  The resulting birational
selfmaps of a fixed cover move a general point densely through
that cover, since their projected graphs dominate \(Y\times Y\).

The cover is non-uniruled, being generically finite over the
non-uniruled \(Y\).  Hanamura's non-uniruled birational-group
theorem gives a smooth projective birational model for which the
reduced birational group is a group scheme, locally of finite type,
and its identity component is an abelian variety
\cite[Theorems~2.1--2.2]{Proj-Hanamura1988};
see also \cite[Proposition~3.7 and Theorems~3.8--3.9]{Proj-Blanc2017}.
After conjugating the maps to that model, shrink the irreducible
parameter variety anew so that their graph closures are flat
and both graph projections remain birational on every fiber.
They then define a morphism to the represented birational scheme;
because the parameter variety is reduced, this morphism factors
through its reduction.  Its connected image lies in a single
component, hence in a translate of the identity component.
Dominance of evaluation makes the corresponding
abelian variety action generically transitive.  Its orbit is an
abelian-variety quotient, so the fixed cover is birational to an
abelian variety.

This is impossible for the original counterexample.  Resolve the
generically finite rational map from that abelian variety to \(X\):
on a smooth model \(U\), ramification gives
\[
 K_U=f^*K_X+R,\qquad R\ge0,
\]
whereas \(K_U\) is an effective divisor exceptional over the
abelian variety.  Let \(T\) be a positive current in \(K_X\).
The current \(f^*T+[R]\) is positive and represents \(K_U\).
Intersecting with powers of the ample class pulled back from the
abelian variety shows that this sum is supported on the exceptional
locus.  Its positive summand \(f^*T\) is therefore supported there
as well, and the support theorem makes \(f^*T\) divisorial.
Push forward this summand alone: \(f_*(f^*T)\) is a positive
divisorial current whose class is \(\deg(f)c_1(K_X)\) by the
projection formula.  Rationality of the numerical class supplies
a rational effective representative, and irregularity zero turns
numerical effectivity into nonvanishing, as in
Lemma~\ref{proj:ex:irregularity}.  This contradicts
\(\kappa(X,K_X)=-\infty\).

All possible \(V\) have now been excluded.  This proves the
Seshadri assertion.  At a very general point the big semiample
contraction is an isomorphism onto its image near that point.
The jet interpretation of the Seshadri constant for the ample
class downstairs gives the stated divisible jet-separating
multiple.  Complete systems agree under the model comparisons,
so this conclusion transfers to the original torus open and to
common resolutions.
\end{proof}
